% Generated by tools/textbook/export_listings.js from tools/concepts/assessments.json: do not edit.
\begin{chapstatus}{Il filo del capitolo}
Il corso segue un solo metodo: giocare, osservare, formalizzare, ritornare. Ogni argomento comincerà come qualcosa visto su PRG32 e finirà come una regola che permette di prevederlo. Prossimo passo: installare il laboratorio.
\par\medskip\textbf{Biglietto d'uscita} (da due a quattro minuti, da soli, su carta):
\begin{enumerate}[nosep,leftmargin=2em]
\item Prima di eseguire un programma vi si chiede di scrivere che cosa vi aspettate di vedere. Quale passo della spirale di apprendimento è questo, e che cosa imparate quando la previsione si rivela sbagliata?
\item Vero o falso? Correggete le affermazioni false. (a) Superare la prova intercorso basta per verbalizzare il voto finale. (b) Nella prova intercorso è ammessa la calcolatrice del telefono. (c) La parte di teoria si può sostenere prima di quella di laboratorio.
\end{enumerate}
\end{chapstatus}
